% GENERATED FILE. Edit canonical lessons, then run npm run generate:books.
% canonical-source-hash: fnv1a64:36f50789e1edb315
% canonical-lessons: KA-C107-bala, KA-C107-eda, KA-C107-nera, KA-C107-dura, KA-C107-mele

\chapter{Right, Left, Straight, Far, Above}
\label{ch:ka-right-left-straight-far-above}
\hlchaptermodality{107}

\begin{chapteropening}
\textbf{By the end of this chapter:} \emph{I can say right, left, straight, far and above.}

Complete the last lesson of chapter 107: i can say right, left, straight, far and above.
\end{chapteropening}

\section[\texorpdfstring{\kn{ಬಲ} (bala)}{ಬಲ (bala)}]{\kn{ಬಲ} (bala) — right (not left)}

\label{lesson:KA-C107-bala}



\begin{quote}
Before the new one: say the Kannada for please phone, then the Kannada for I'll see you.
\end{quote}

\subsection*{You'll want to know: \kn{ಬಲ}}
\textbf{\kn{ಬಲ}} — \emph{bala} — \textquotedblleft{}right (not left)\textquotedblright{}.

\begin{grammarlens}[title={What you've built}]
One more word for this chapter.
\end{grammarlens}

\subsection*{Your turn}
\begin{itemize}
\raggedright
  \item \emph{Say it:} \emph{bala}
  \item \emph{Say it:} \emph{bala}, once more
  \item \emph{Say it:} say \emph{bala}
  \item \emph{Recall:} say the Kannada for please phone, then the Kannada for I'll see you, then say \emph{bala} again
\end{itemize}

\begin{tcolorbox}[breakable,colback=teal!4,colframe=teal!35!black,title={Before you move on}]
What does \kn{ಬಲ} mean? (\textquotedblleft{}Right (not left)\textquotedblright{}.) Say it once more.
\end{tcolorbox}
\section[\texorpdfstring{\kn{ಎಡ} (eḍa)}{ಎಡ (eḍa)}]{\kn{ಎಡ} (eḍa) — left (not right)}

\label{lesson:KA-C107-eda}



\begin{quote}
Before the new one: say the Kannada for I'll see you, then the Kannada for right.
\end{quote}

\subsection*{You'll want to know: \kn{ಎಡ}}
\textbf{\kn{ಎಡ}} — \emph{eḍa} — \textquotedblleft{}left (not right)\textquotedblright{}.

\begin{grammarlens}[title={What you've built}]
One more word for this chapter.
\end{grammarlens}

\subsection*{Your turn}
\begin{itemize}
\raggedright
  \item \emph{Say it:} \emph{eḍa}
  \item \emph{Say it:} \emph{eḍa}, once more
  \item \emph{Say it:} say \emph{eḍa}
  \item \emph{Recall:} say the Kannada for I'll see you, then the Kannada for right, then say \emph{eḍa} again
\end{itemize}

\begin{tcolorbox}[breakable,colback=teal!4,colframe=teal!35!black,title={Before you move on}]
What does \kn{ಎಡ} mean? (\textquotedblleft{}Left (not right)\textquotedblright{}.) Say it once more.
\end{tcolorbox}
\section[\texorpdfstring{\kn{ನೇರ} (nēra)}{ನೇರ (nēra)}]{\kn{ನೇರ} (nēra) — straight}

\label{lesson:KA-C107-nera}



\begin{quote}
Before the new one: say the Kannada for right, then the Kannada for left.
\end{quote}

\subsection*{You'll want to know: \kn{ನೇರ}}
\textbf{\kn{ನೇರ}} — \emph{nēra} — \textquotedblleft{}straight\textquotedblright{}.

\begin{grammarlens}[title={What you've built}]
One more word for this chapter.
\end{grammarlens}

\subsection*{Your turn}
\begin{itemize}
\raggedright
  \item \emph{Say it:} \emph{nēra}
  \item \emph{Say it:} \emph{nēra}, once more
  \item \emph{Say it:} say \emph{nēra}
  \item \emph{Recall:} say the Kannada for right, then the Kannada for left, then say \emph{nēra} again
\end{itemize}

\begin{tcolorbox}[breakable,colback=teal!4,colframe=teal!35!black,title={Before you move on}]
What does \kn{ನೇರ} mean? (\textquotedblleft{}Straight\textquotedblright{}.) Say it once more.
\end{tcolorbox}
\section[\texorpdfstring{\kn{ದೂರ} (dūra)}{ದೂರ (dūra)}]{\kn{ದೂರ} (dūra) — far; distance}

\label{lesson:KA-C107-dura}



\begin{quote}
Before the new one: say the Kannada for left, then the Kannada for straight.
\end{quote}

\subsection*{You'll want to know: \kn{ದೂರ}}
\textbf{\kn{ದೂರ}} — \emph{dūra} — \textquotedblleft{}far; distance\textquotedblright{}.

\begin{grammarlens}[title={What you've built}]
One more word for this chapter.
\end{grammarlens}

\subsection*{Your turn}
\begin{itemize}
\raggedright
  \item \emph{Say it:} \emph{dūra}
  \item \emph{Say it:} \emph{dūra}, once more
  \item \emph{Say it:} say \emph{dūra}
  \item \emph{Recall:} say the Kannada for left, then the Kannada for straight, then say \emph{dūra} again
\end{itemize}

\begin{tcolorbox}[breakable,colback=teal!4,colframe=teal!35!black,title={Before you move on}]
What does \kn{ದೂರ} mean? (\textquotedblleft{}Far; distance\textquotedblright{}.) Say it once more.
\end{tcolorbox}
\section[\texorpdfstring{\kn{ಮೇಲೆ} (mēle)}{ಮೇಲೆ (mēle)}]{\kn{ಮೇಲೆ} (mēle) — above, on top}

\label{lesson:KA-C107-mele}



\begin{quote}
Before the new one: say the Kannada for straight, then the Kannada for far; distance.
\end{quote}

\subsection*{You'll want to know: \kn{ಮೇಲೆ}}
\textbf{\kn{ಮೇಲೆ}} — \emph{mēle} — \textquotedblleft{}above, on top\textquotedblright{}.

\begin{grammarlens}[title={What you've built}]
One more word for this chapter.
\end{grammarlens}

\subsection*{Your turn}
\begin{itemize}
\raggedright
  \item \emph{Say it:} \emph{mēle}
  \item \emph{Say it:} \emph{mēle}, once more
  \item \emph{Say it:} say \emph{mēle}
  \item \emph{Recall:} say the Kannada for straight, then the Kannada for far; distance, then say \emph{mēle} again
\end{itemize}

\begin{tcolorbox}[breakable,colback=teal!4,colframe=teal!35!black,title={Before you move on}]
What does \kn{ಮೇಲೆ} mean? (\textquotedblleft{}Above, on top\textquotedblright{}.) Say it once more.
\end{tcolorbox}
